\chapter{Elementary representations of groups}
\label{chap:M07.17}
The basic constructions in this chapter apply to any group. Finiteness of
the group and restrictions on the characteristic are introduced only in
the averaging theorem. Scalar endomorphism and dimension-one conclusions
have their own algebraic-closedness and dimension hypotheses.

\KEntry{K07.17.definition}{Linear and matrix representations}
\begin{definition}
For a field $K$, a group $G$ and a $K$-vector space $V$, a representation
is a homomorphism $\rho:G\to\GL_K(V)$, where the product of linear
automorphisms is composition. We write $gv=\rho(g)v$.
\end{definition}
Equivalently, one may give endomorphisms $\rho(g)\in\End_K(V)$
satisfying $\rho(1)=I$ and $\rho(gh)=\rho(g)\rho(h)$.
Indeed $\rho(g)\rho(g^{-1})=\rho(g^{-1})\rho(g)=I$, so every
endomorphism so specified is invertible with inverse $\rho(g^{-1})$.
Conversely an automorphism-valued homomorphism satisfies these identities
after forgetting the inverse data. These constructions are inverse.

For a finite basis $b=(b_1,\ldots,b_n)$, let $A_g$ be the matrix of
$\rho(g)$ in $b$. The rule for matrices of compositions gives
$A_{gh}=A_gA_h$ and $A_1=I_n$, and $A_{g^{-1}}$ is the inverse of
$A_g$. Conversely a homomorphism $G\to\GL_n(K)$ acts on $K^n$ and
can be transported to $V$ by the coordinate isomorphism. Thus the
chosen basis provides the stated correspondence, including $n=0$.

\KEntry{K07.17.examples}{Trivial, permutation, regular and sign representations}
The trivial representation has $\rho(g)=I$ for every $g$; its
homomorphism laws follow from $I^2=I$. If $X$ is a $G$-set, let
$K^{(X)}$ be the space of finite-support formal sums $\sum_xa_xe_x$.
Set $\rho(g)e_x=e_{gx}$ and extend linearly. A finite-support vector
has a finite-support image. On every basis vector,
$\rho(g)\rho(h)e_x=e_{g(hx)}=e_{(gh)x}$ and $\rho(1)e_x=e_x$;
equality on the basis proves the representation laws. For finite $X$
the dimension is $|X|$. For $X=G$ with left translation this is the
regular representation. If $g$ acts trivially in the regular
representation, then $e_g=\rho(g)e_1=e_1$, so $g=1$; it is faithful.
For a general $G$-set, an element fixes every basis vector exactly
when it fixes every point, so the kernel is the action kernel.

The sign homomorphism of the symmetric group gives a one-dimensional
representation $\rho(\sigma)v=\sgn(\sigma)v$, with $\pm1$ mapped
to $K$. Multiplicativity of sign is proved in
\ref{K07.02.transpositions_parity}; scalar multiplication then proves
the representation law. In characteristic two $-1=1$, so this
representation is trivial. No irreducibility claim follows merely
from the existence of a permutation basis.

\KEntry{K07.17.subrepresentations_quotients}{Invariant subspaces and quotients}
A subspace $U\le V$ is invariant if $\rho(g)U\subseteq U$ for
every $g$. This is in fact equality: for $u\in U$,
$u=\rho(g)\rho(g^{-1})u$ lies in $\rho(g)U$. Restriction therefore
defines an automorphism of $U$, and restriction respects products and
identity. On the quotient set of additive cosets put
\[
 g(v+U)=\rho(g)v+U.
\]
If $v-w\in U$, then $\rho(g)v-\rho(g)w\in U$; this proves
representative independence. Additivity and compatibility with scalars
descend from $\rho(g)$. Identity and composition can be checked on
$v+U$ using those laws in $V$, and the action of $g^{-1}$ is inverse.
Thus $V/U$ has a quotient representation. The inclusion of $U$ and
the quotient map are equivariant. The condition is invariance of a
subspace, not normality of a subgroup. For monoid actions the equality
$\rho(g)U=U$ need not follow from inclusion.

\KEntry{K07.17.intertwiners_equivalence}{Intertwiners and equivalence}
For representations $\rho$ on $V$ and $\sigma$ on $W$, an
intertwiner is a linear map $T:V\to W$ satisfying
\[
 T\rho(g)=\sigma(g)T\quad(g\in G).
\]
Intertwiners are closed under addition and scalar multiplication, by
distributing both sides of this equality. Identity maps intertwine,
and the composite of two intertwiners intertwines, by substituting
the two equations. If $T$ is bijective, the equation multiplied by
$T^{-1}$ on both sides gives
$T^{-1}\sigma(g)=\rho(g)T^{-1}$. Consequently invertible
intertwiners define a reflexive, symmetric and transitive equivalence
relation on representations.

Suppose $V,W$ have finite bases and equal dimension $n$; write their
matrices as $A_g,B_g$ and the matrix of $T$ as $P$. Equality of linear
maps is equivalent to equality of their matrices, so the intertwining
equations are $PA_g=B_gP$. The map $T$ is invertible exactly when
$P$ is invertible, in which case
$B_g=PA_gP^{-1}$ for every $g$, with the \emph{same} $P$.
This also proves the converse by turning $P$ into a coordinate map.
Representations of different finite dimensions cannot be equivalent,
by the basis theorem. Choosing a different conjugating matrix for
each group element does not establish an equivalence.

\KEntry{K07.17.kernels_images}{Kernels, images and the first isomorphism}
\begin{proposition}
If $T$ intertwines $\rho$ and $\sigma$, its kernel and image are
invariant, and the linear isomorphism
$\overline T:V/\ker T\to\im T$ is equivariant.
\end{proposition}
\begin{proof}
If $Tv=0$, then $T\rho(g)v=\sigma(g)Tv=0$.
If $w=Tv$, then $\sigma(g)w=T\rho(g)v\in\im T$.
These prove invariance. The linear first isomorphism theorem
identifies $\overline T(v+\ker T)=Tv$ as a well-defined bijective
linear map. For its additional equivariance property compute
\[
 \overline T(g(v+\ker T))=T\rho(g)v
   =\sigma(g)Tv=g\overline T(v+\ker T).
\]
Every quotient class has such a representative, so this proves the
equation on the whole quotient.
\end{proof}

\KEntry{K07.17.direct_sums}{Direct sums and a common block diagonal form}
For representations on $V,W$, define
$g(v,w)=(\rho(g)v,\sigma(g)w)$. Linearity, identity and
composition follow componentwise. The two inclusions and projections
intertwine. For an internal decomposition into invariant subspaces
$E=U\oplus W$, the addition map $U\oplus W\to E$ is an
equivariant linear isomorphism, because
$g(u+w)=gu+gw$ and uniqueness of direct-sum coordinates gives the
inverse. If bases of $U,W$ are finite, their concatenation is a basis
of $E$ in which \emph{all} action matrices are block diagonal.

Conversely, a common basis in which every $A_g$ has the same block
partition determines the spans of the two basis blocks. Multiplication
by each block diagonal matrix preserves both spans. They have zero
intersection and span $E$, proving the invariant internal direct sum.
An invariant subspace alone gives only an upper triangular block
form after extending its basis: the other block need not be invariant.
For an infinite family, direct sum denotes finite support, and no
equality with the full product is asserted.

\KEntry{K07.17.irreducible_definition}{Irreducibility and semisimplicity}
A representation on $V$ is irreducible when $V\ne0$ and its only
invariant subspaces are $0,V$. It is semisimple when every invariant
subspace has an invariant complement. A zero representation space is
semisimple and is not irreducible.

\begin{theorem}\label{thm:semisimple-finite-sum}
In finite dimension, a representation is semisimple if and only if it
is a finite direct sum of irreducible representations. No finiteness
of $G$ or characteristic restriction is needed for this equivalence.
\end{theorem}
\begin{proof}
Suppose first that $V$ is semisimple. If $V=0$, use the empty sum.
Otherwise choose a nonzero invariant subspace $S$ of least positive
dimension. Any proper nonzero invariant subspace of $S$ would have
smaller positive dimension, so $S$ is irreducible. Choose an invariant
complement $T$. It too is semisimple: for $A\le T$ invariant, choose
$V=A\oplus B$ with $B$ invariant. Every $t\in T$ is $a+b$ with
$a\in A\subset T$, hence $b=t-a\in T\cap B$. Thus
$T=A\oplus(T\cap B)$. Since $\dim T<\dim V$, dimension induction
decomposes $T$ and then $V$ as required.

For the converse induct on the number of irreducible summands. The
empty sum has the asserted property. Write $V=S\oplus T$ with $S$
irreducible and $T$ a sum of fewer irreducibles, hence semisimple by
induction. Let $A\le V$ be invariant. Its intersection with $S$ is
either $S$ or zero. If $S\subset A$, projecting onto $T$ shows
$A=S\oplus(A\cap T)$: subtracting the $S$-coordinate leaves a
vector of $A\cap T$. Choose an invariant complement $C$ to $A\cap T$
in $T$. Then $V=A\oplus C$.

If $A\cap S=0$, the projection $q:V\to T$ restricts to an
equivariant isomorphism $A\to B=q(A)$. Since $T$ is semisimple,
$T=B\oplus C$ for an invariant $C$. Relative to $S\oplus T$, the
inverse of $q|_A$ writes $A$ as the graph
$\{(f(b),b):b\in B\}$ of a linear map $f:B\to S$.
For $(s,b+c)\in S\oplus(B\oplus C)$ the decomposition
\[
 (s,b+c)=(f(b),b)+(s-f(b),c)
\]
shows $V=A+(S\oplus C)$. Intersection is zero because its
$T$-coordinate would lie in $B\cap C=0$ and its $A$-vector would
then lie in $A\cap S=0$. Both $S$ and $C$ are invariant, so
$S\oplus C$ is the required invariant complement.
\end{proof}
Irreducibility is stronger than indecomposability: the modular example
below has no invariant complementary line but does have an invariant
line. Over $\R$, let a generator of $C_4$ act on $\R^2$ by
$J=\left(\begin{smallmatrix}0&-1\\1&0\end{smallmatrix}\right)$.
An invariant real line would give a real eigenvalue $a$ with $a^2=-1$,
which is impossible. Over $\C$ the eigenvalues $i,-i$ give two
invariant eigenlines. Thus extending scalars can change irreducibility.

\KEntry{K07.17.schur_nonscalar}{Schur's zero-or-isomorphism lemma}
\begin{theorem}\label{thm:schur-general}
A nonzero intertwiner between irreducible representations over any
field is an equivalence. In particular $\End_G(V)$ is a division
algebra for an irreducible $V$, without a finite-dimensional hypothesis.
\end{theorem}
\begin{proof}
The kernel and image are invariant by the earlier proposition. The
kernel cannot be all of $V$, because $T\ne0$; it is therefore zero.
The image is nonzero, so irreducibility of the target makes it the
whole target. The map is thus bijective; its linear inverse
intertwines by the inverse computation already proved. Equivariant
endomorphisms form a subalgebra of $\End_K(V)$, and every nonzero
element now has an inverse in that subalgebra. Its identity is
nonzero since $V\ne0$. This is precisely a division algebra.
\end{proof}
The real rotation example supplies a nonscalar equivariant
endomorphism $J$ with $J^2=-I$, since $J$ commutes with all powers
of $J$. Hence a scalar conclusion would be false over an arbitrary
field.

\KEntry{K07.17.schur_scalar}{Scalar Schur lemma}
\begin{theorem}\label{thm:schur-scalar}
If $K$ is algebraically closed and $V$ is nonzero, finite-dimensional
and irreducible, each $T\in\End_G(V)$ is a unique scalar multiple
of $I$. The scalar map $K\to\End_G(V)$ is an algebra isomorphism.
\end{theorem}
\begin{proof}
The characteristic polynomial has positive degree and a root
$\lambda\in K$, so $T-\lambda I$ has nonzero kernel by the
linear-algebra foundation. It is an intertwiner because both $T$
and $\lambda I$ commute with the action. If it were nonzero,
Theorem~\ref{thm:schur-general} would make it invertible, contradicting
its nonzero kernel. Thus $T=\lambda I$. If also $T=\mu I$, choose
$v\ne0$; then $(\lambda-\mu)v=0$ forces $\lambda=\mu$ in a
vector space over a field. Scalar multiplication preserves addition,
product and identity and lands in the commutant. The preceding
arguments prove surjectivity and injectivity, respectively. Its
$K$-linear dimension is therefore one.
\end{proof}

\KEntry{K07.17.maschke_projection}{Maschke averaging and invariant complements}
\begin{theorem}\label{thm:maschke}
Let $G$ be finite and $K$ a field such that $m=|G|\cdot1_K\ne0$.
Every invariant subspace $U$ of any $K$-representation $V$ has an
invariant complement, even when $V$ is infinite-dimensional.
\end{theorem}
\begin{proof}
Extend a basis of $U$ to a basis of $V$ and let $E:V\to V$ be the
projection onto $U$ that kills the added basis vectors. Thus
$E(V)\subset U$ and $E|_U=I_U$. Define
\[
 P=m^{-1}\sum_{g\in G}\rho(g)E\rho(g^{-1}).
\]
Every summand has image in $U$, since $U$ is invariant. For $u\in U$,
$\rho(g^{-1})u\in U$, so each summand sends $u$ to $u$.
Consequently $P(V)\subset U$ and $P|_U=I_U$, whence $P^2=P$.
For $h\in G$, substitution $g'=hg$ is a bijection of $G$ and gives
\[
 \rho(h)P\rho(h^{-1})
  =m^{-1}\sum_{g\in G}\rho(hg)E\rho((hg)^{-1})=P.
\]
Therefore $P$ commutes with every $\rho(h)$ and $\ker P$ is
invariant. Finally every $v$ is
$P(v)+(v-P(v))$, with first term in $U$ and second in $\ker P$.
An element of $U\cap\ker P$ equals its own image under $P$ and
also has image zero, so it is zero. Thus $V=U\oplus\ker P$.
\end{proof}
The nonzero scalar condition holds in characteristic zero. In positive
characteristic $p$, it is exactly $p\nmid|G|$; division by $|G|$ is
not meaningful when this condition fails.

\KEntry{K07.17.complete_reducibility}{Finite-dimensional complete reducibility}
Under the hypotheses of Theorem~\ref{thm:maschke}, a
finite-dimensional representation is a finite direct sum of
irreducibles. Indeed Maschke gives the invariant-complement
definition of semisimplicity, and
Theorem~\ref{thm:semisimple-finite-sum} proves the required finite
decomposition, including the empty sum for $V=0$. The equivalence of
that internal sum with an external sum is the addition isomorphism
proved in \ref{K07.17.direct_sums}. This argument requires neither
algebraic closedness nor that all irreducible summands be one-dimensional.

\KEntry{K07.17.modular_counterexample}{Failure in the modular characteristic}
Let $p$ be prime and work over $K=\F_p$. Put
$N=\left(\begin{smallmatrix}0&1\\0&0\end{smallmatrix}\right)$
and $J=I+N$. Since $N^2=0$, induction on $a\ge0$ gives
$(I+N)^a=I+aN$: the inductive step multiplies out and the $N^2$
term vanishes. In particular $J^p=I$. Hence the rule sending the
generator of $C_p=\Z/p\Z$ to $J$ is well-defined: changing an
integer exponent by a multiple of $p$ leaves its power unchanged.
It is a representation because exponents add when powers multiply.
The line $U=Ke_1$ is invariant and fixed pointwise.

Suppose $W$ is a complementary invariant line. Its nonzero vector
has nonzero second coordinate, since otherwise $W=U$; rescale it
to $w=ae_1+e_2$. Then $Jw=(a+1)e_1+e_2$. Invariance makes
$Jw=cw$ for some $c\in K$. Comparing second coordinates gives
$c=1$; comparing first coordinates then gives $a+1=a$, impossible
in a field. Thus no invariant complement exists. The argument works
for $p=2$ as well. The same calculation shows the only invariant
line is $U$, so a decomposition into two nonzero invariant summands
is impossible, although the representation is reducible.

\KEntry{K07.17.abelian_irreducible_dimension}{Irreducibles of an abelian group}
Suppose $G$ is abelian, $K$ algebraically closed, and $V$ a nonzero
finite-dimensional irreducible representation. For each $g$, the
operator $\rho(g)$ commutes with every $\rho(h)$, so scalar Schur
gives $\rho(g)=\lambda_gI$. Choose $v\ne0$. The line $Kv$ is
invariant under all these scalar operators; irreducibility forces
$Kv=V$. Thus $\dim_KV=1$. No finiteness of $G$ was used. In
characteristic dividing a finite group order this conclusion about
irreducibles remains valid, while arbitrary representations need
not split, as the preceding example shows. The real rotation
example proves the algebraic-closedness hypothesis cannot simply
be omitted.

\KEntry{K07.17.characters}{Characters and conjugacy invariance}
For finite-dimensional $V$ define the character
$\chi_\rho(g)=\Tr_K(\rho(g))$. The trace is basis independent.
For completeness, in finite matrices
\[
 \Tr(AB)=\sum_i\sum_j a_{ij}b_{ji}
        =\sum_j\sum_i b_{ji}a_{ij}=\Tr(BA),
\]
where commuting scalar entries and interchanging finite sums justify
the equality. Hence $\Tr(PAP^{-1})=\Tr(A)$.
Applying this with $P=\rho(h)$ proves
$\chi_\rho(hgh^{-1})=\chi_\rho(g)$. An equivalence has action
matrices conjugate by a single $P$, so it preserves characters as
well. Since $\rho(1)=I$, $\chi_\rho(1)=(\dim_KV)\cdot1_K$.
In positive characteristic this value need not recover the natural
number dimension. For example, over $\F_p$ the zero space and a
$p$-dimensional trivial representation of the trivial group both
have zero character and are not equivalent. Thus no arbitrary-field
character-classification theorem is asserted here.

\KEntry{K07.17.character_direct_sum}{Trace of a direct sum and permutation characters}
In concatenated finite bases, the action of $g$ on $V\oplus W$
has blocks $\rho(g),\sigma(g)$. Summing its diagonal entries
therefore proves
\[
 \chi_{\rho\oplus\sigma}(g)=\chi_\rho(g)+\chi_\sigma(g).
\]
This identity uses no averaging or characteristic hypothesis. For
the representation associated to a finite $G$-set $X$, the matrix
column of $e_x$ has its single $1$ in row $gx$. Its diagonal entry
is $1$ exactly when $gx=x$, and otherwise is zero. Consequently
\[
 \chi(g)=|\{x\in X:gx=x\}|\cdot1_K.
\]
This is a count of fixed \emph{basis points}, interpreted in $K$.
It is not a formula equating trace and dimension of a fixed subspace
for an arbitrary linear operator.
